%%%%%%%%%%%%%%%%
%%% deut 8 %%%
%%%%%%%%%%%%%%%%

\Chapter{8}

\Verse{1}
{
	\hDtrA{כּל הַמִּצְוָה אֲשֶׁר אָנֹכִי מְצַוְּךָ הַיּוֹם תִּשְׁמְרוּן לַעֲשׂוֹת לְמַעַן תִּחְיוּן וּרְבִיתֶם וּבָאתֶם וִירִשְׁתֶּם אֶת הָאָרֶץ אֲשֶׁר נִשְׁבַּע יְהֹוָה לַאֲבֹתֵיכֶם׃}
}
{
	\eDtrA{
		“All of the commandment that I command you today you
		shall be watchful to do, so that you’ll live, and you’ll
		multiply, and you’ll come and take possession of the land
		that YHWH swore to your fathers.
	}
}
{
}

\Verse{2}
{
	\hDtrA{וְזָכַרְתָּ אֶת כּל הַדֶּרֶךְ אֲשֶׁר הוֹלִיכְךָ יְהֹוָה אֱלֹהֶיךָ זֶה אַרְבָּעִים שָׁנָה בַּמִּדְבָּר לְמַעַן עַנֹּתְךָ לְנַסֹּתְךָ לָדַעַת אֶת אֲשֶׁר בִּלְבָבְךָ הֲתִשְׁמֹר מִצְוֺתָו אִם לֹא׃}
}
{
	\eDtrA{
		And you shall remember all the way that YHWH, your God,
		had you go these forty years in the wilderness in order to
		degrade you, to test you, to know what was in your heart:
		would you observe His commandments or not.
	}
}
{
%	\footnote{
%		to degrade you, to test you, to know. This sequence of verbs is powerful
%		because of the ways that they are used before this. The word “degrade”
%		is used to describe what the Egyptians did to the Israelites as slaves
%		(Exod 1:11). (It is also used for what Shechem did to Dinah, and for
%		what the people of Israel are supposed to do to themselves on Yom
%		Kippur.) It means to bring down, to humble, to lower someone. The
%		difference is that the Egyptians do it to weaken the Israelites, whereas
%		God does it to test their strength. The second word here, “test,” is the
%		word for what God does in commanding Abraham to sacrifice Isaac (Gen
%		22:1). The result of that test is that God says, “Now I know that you
%		fear God” (22:12). And so the third word here, likewise, is “know,”
%		conveying that the purpose of a divine test is still to establish what
%		is in a human’s heart. Moses thus now tells the people that the reason
%		for many of the tribulations that they experienced for forty years was
%		to test them. We, the readers, were informed of this at the beginning of
%		the account of the years in the wilderness, where testing is mentioned
%		twice (Exod 15:25; 16:4), but this is the first time that Moses reveals
%		it to the people.
%	}
}

\Verse{3}
{
	\hDtrA{וַיְעַנְּךָ וַיַּרְעִבֶךָ וַיַּאֲכִלְךָ אֶת הַמָּן אֲשֶׁר לֹא יָדַעְתָּ וְלֹא יָדְעוּן אֲבֹתֶיךָ לְמַעַן הוֹדִיעֲךָ כִּי לֹא עַל הַלֶּחֶם לְבַדּוֹ יִחְיֶה הָאָדָם כִּי עַל כּל מוֹצָא פִי יְהֹוָה יִחְיֶה הָאָדָם׃}
}
{
	\eDtrA{
		So He degraded you and made you hunger and then fed you
		the manna, which you had not known and your fathers had
		not known, to let you know that a human doesn’t live by
		bread alone, but a human lives by every product of YHWH’s
		mouth.
	}
}
{
}

\Verse{4}
{
	\hDtrA{שִׂמְלָתְךָ לֹא בָלְתָה מֵעָלֶיךָ וְרַגְלְךָ לֹא בָצֵקָה זֶה אַרְבָּעִים שָׁנָה׃}
}
{
	\eDtrA{
		Your garment didn’t wear out on you, and your foot didn’t
		swell these forty years.
	}
}
{
}

\Verse{5}
{
	\hDtrA{וְיָדַעְתָּ עִם לְבָבֶךָ כִּי כַּאֲשֶׁר יְיַסֵּר אִישׁ אֶת בְּנוֹ יְהֹוָה אֱלֹהֶיךָ מְיַסְּרֶךָּ׃}
}
{
	\eDtrA{
		So you shall know in your heart that, the way a man
		disciplines his child, so YHWH, your God, disciplines you.
	}
}
{
}

\Verse{6}
{
	\hDtrA{וְשָׁמַרְתָּ אֶת מִצְוֺת יְהֹוָה אֱלֹהֶיךָ לָלֶכֶת בִּדְרָכָיו וּלְיִרְאָה אֹתוֹ׃}
}
{
	\eDtrA{
		And you shall observe the commandments of YHWH, your God,
		to go in His ways and to fear Him.
	}
}
{
%	\footnote{
%		fear. Some may find the idea of being afraid of God to be strange or
%		unattractive. But the fact remains that to experience the power of the
%		creator of the universe is fearful. Thus Moses hides his face at the
%		burning bush because he is afraid to look at God. Thus the people are
%		terrified when they hear the divine voice at Sinai. Thus Nadab and Abihu
%		die because of a misstep in closeness to the holy. When one appreciates
%		what is at stake, one can understand why the idea of fear of
%		God—literally, profoundly—is so significant here.
%	}
}

\Verse{7}
{
	\hDtrA{כִּי יְהֹוָה אֱלֹהֶיךָ מְבִיאֲךָ אֶל אֶרֶץ טוֹבָה אֶרֶץ נַחֲלֵי מָיִם עֲיָנֹת וּתְהֹמֹת יֹצְאִים בַּבִּקְעָה וּבָהָר׃}
}
{
	\eDtrA{
		Because YHWH, your God, is bringing you to a good land, a
		land of wadis of water, springs, and deeps coming out in
		valleys and in hills,
	}
}
{
}

\Verse{8}
{
	\hDtrA{אֶרֶץ חִטָּה וּשְׂעֹרָה וְגֶפֶן וּתְאֵנָה וְרִמּוֹן אֶרֶץ זֵית שֶׁמֶן וּדְבָשׁ׃}
}
{
	\eDtrA{
		a land of wheat and barley and vine and fig and
		pomegranate, a land of olive, oil, and honey,
	}
}
{
}

\Verse{9}
{
	\hDtrA{אֶרֶץ אֲשֶׁר לֹא בְמִסְכֵּנֻת תֹּאכַל בָּהּ לֶחֶם לֹא תֶחְסַר כֹּל בָּהּ אֶרֶץ אֲשֶׁר אֲבָנֶיהָ בַרְזֶל וּמֵהֲרָרֶיהָ תַּחְצֹב נְחֹשֶׁת׃}
}
{
	\eDtrA{
		a land in which you won’t eat bread in scarcity, in which
		you won’t lack anything, a land whose stones are iron and
		from whose hills you’ll hew bronze.
	}
}
{
%	\footnote{
%		you’ll hew bronze. Copper and tin.
%	}
}

\Verse{10}
{
	\hDtrA{וְאָכַלְתָּ וְשָׂבָעְתָּ וּבֵרַכְתָּ אֶת יְהֹוָה אֱלֹהֶיךָ עַל הָאָרֶץ הַטֹּבָה אֲשֶׁר נָתַן לָךְ׃}
}
{
	\eDtrA{
		And you’ll eat and be full and bless YHWH, your God, for
		the good land that He has given you.
	}
}
{
%	\footnote{
%		eat and be full and bless. The Jewish practice of saying grace after
%		meals, birkt hammzn, is associated with this verse.
%	}
}

\Verse{11}
{
	\hDtrA{הִשָּׁמֶר לְךָ פֶּן תִּשְׁכַּח אֶת יְהֹוָה אֱלֹהֶיךָ לְבִלְתִּי שְׁמֹר מִצְוֺתָיו וּמִשְׁפָּטָיו וְחֻקֹּתָיו אֲשֶׁר אָנֹכִי מְצַוְּךָ הַיּוֹם׃}
}
{
	\eDtrA{
		“Watch yourself in case you’ll forget YHWH, your God, so
		as not to observe His commandments and His judgments and
		His laws that I command you today,
	}
}
{
}

\Verse{12}
{
	\hDtrA{פֶּן תֹּאכַל וְשָׂבָעְתָּ וּבָתִּים טֹבִים תִּבְנֶה וְיָשָׁבְתָּ׃}
}
{
	\eDtrA{
		in case you’ll eat, and you’ll be full, and you’ll build
		good houses, and you’ll live there,
	}
}
{
}

\Verse{13}
{
	\hDtrA{וּבְקָרְךָ וְצֹאנְךָ יִרְבְּיֻן וְכֶסֶף וְזָהָב יִרְבֶּה לָּךְ וְכֹל אֲשֶׁר לְךָ יִרְבֶּה׃}
}
{
	\eDtrA{
		and your oxen and your flock will multiply, and silver and
		gold will multiply for you, and everything that you’ll
		have will multiply,
	}
}
{
}

\Verse{14}
{
	\hDtrA{וְרָם לְבָבֶךָ וְשָׁכַחְתָּ אֶת יְהֹוָה אֱלֹהֶיךָ הַמּוֹצִיאֲךָ מֵאֶרֶץ מִצְרַיִם מִבֵּית עֲבָדִים׃}
}
{
	\eDtrA{
		and your heart will be lifted, and you’ll forget YHWH,
		your God, who brought you out from the land of Egypt, from
		a house of slaves,
	}
}
{
}

\Verse{15}
{
	\hDtrA{הַמּוֹלִיכְךָ בַּמִּדְבָּר הַגָּדֹל וְהַנּוֹרָא נָחָשׁ שָׂרָף וְעַקְרָב וְצִמָּאוֹן אֲשֶׁר אֵין מָיִם הַמּוֹצִיא לְךָ מַיִם מִצּוּר הַחַלָּמִישׁ׃}
}
{
	\eDtrA{
		who led you in the big and fearful wilderness of fiery
		snake and scorpion and thirst, where there isn’t water,
		who brought out water for you from the flint rock,
	}
}
{
}

\Verse{16}
{
	\hDtrA{הַמַּאֲכִלְךָ מָן בַּמִּדְבָּר אֲשֶׁר לֹא יָדְעוּן אֲבֹתֶיךָ לְמַעַן עַנֹּתְךָ וּלְמַעַן נַסֹּתֶךָ לְהֵיטִבְךָ בְּאַחֲרִיתֶךָ׃}
}
{
	\eDtrA{
		who fed you manna in the wilderness, which your fathers
		had not known, in order to degrade you and in order to
		test you to be good to you in your future;
	}
}
{
}

\Verse{17}
{
	\hDtrA{וְאָמַרְתָּ בִּלְבָבֶךָ כֹּחִי וְעֹצֶם יָדִי עָשָׂה לִי אֶת הַחַיִל הַזֶּה׃}
}
{
	\eDtrA{
		and you’ll say in your heart: ‘My power and my hand’s
		strength made this wealth for me.’
	}
}
{
%	\footnote{
%		My power. The entire passage from v. 11 to v. 17 is one sentence. Moses
%		uses this extraordinarily long sentence to convey the weight of the
%		situation against which he is warning. So much will go well that the
%		people can be drawn to forget the God who has done so much for them, so
%		they think that they have achieved it all on their own. And then Moses
%		answers in a simple sentence of ordinary length: You shall remember that
%		He is the one who gives …
%	}
}

\Verse{18}
{
	\hDtrA{וְזָכַרְתָּ אֶת יְהֹוָה אֱלֹהֶיךָ כִּי הוּא הַנֹּתֵן לְךָ כֹּחַ לַעֲשׂוֹת חָיִל לְמַעַן הָקִים אֶת בְּרִיתוֹ אֲשֶׁר נִשְׁבַּע לַאֲבֹתֶיךָ כַּיּוֹם הַזֶּה׃}
}
{
	\eDtrA{
		Then you shall remember YHWH, your God, because He is the
		one who gave you power to make wealth so as to uphold His
		covenant that He swore to your fathers—as it is this day.
	}
}
{
}

\Verse{19}
{
	\hDtrB{וְהָיָה אִם שָׁכֹחַ תִּשְׁכַּח אֶת יְהֹוָה אֱלֹהֶיךָ וְהָלַכְתָּ אַחֲרֵי אֱלֹהִים אֲחֵרִים וַעֲבַדְתָּם וְהִשְׁתַּחֲוִיתָ לָהֶם הַעִדֹתִי בָכֶם הַיּוֹם כִּי אָבֹד תֹּאבֵדוּן׃}
}
{
	\eDtrB{
		“And it will be, if you’ll forget YHWH, your God, and
		you’ll go after other gods and serve them and bow to them,
		I call witness regarding you today that you’ll perish!
	}
}
{
}

\Verse{20}
{
	\hDtrB{כַּגּוֹיִם אֲשֶׁר יְהֹוָה מַאֲבִיד מִפְּנֵיכֶם כֵּן תֹּאבֵדוּן עֵקֶב לֹא תִשְׁמְעוּן בְּקוֹל יְהֹוָה אֱלֹהֵיכֶם׃}
}
{
	\eDtrB{
		Like the nations that YHWH is making to perish in front of
		you, so will you perish, because you wouldn’t listen to
		the voice of YHWH, your God.
	}
}
{
}
